% Generated draft: requires mathematical and source review.
\documentclass[11pt]{article}
\usepackage[margin=25mm]{geometry}
\usepackage{fontspec}
\setmainfont{Latin Modern Roman}
\usepackage{amsmath,amssymb,mathtools,unicode-math}
\setmathfont{Latin Modern Math}
\usepackage{enumitem,needspace,xurl,hyperref}
\setlength{\parindent}{0pt}
\setlength{\parskip}{5pt}
\setlength{\emergencystretch}{3em}
\newcommand{\N}{\mathbb{N}}
\newcommand{\Z}{\mathbb{Z}}
\newcommand{\Q}{\mathbb{Q}}
\newcommand{\R}{\mathbb{R}}
\newcommand{\C}{\mathbb{C}}
\newcommand{\F}{\mathbb{F}}
\newcommand{\A}{\mathbb{A}}
\newcommand{\PP}{\mathbb{P}}
\newcommand{\E}{\mathbb{E}}
\newcommand{\eps}{\varepsilon}
\newcommand{\dd}{\,\mathrm{d}}
\begin{document}
\section*{UnsolvedMath ID: 30001349}
\textbf{Source problem code:} catalogue entry
\subsection{Definitions and mathematical statement}
For the following interpretation, let \(n\geq 1\) and let \(f\in\Q[x_1,\ldots,x_n]\) be a nonzero polynomial. These domain conventions are not explicit in the supplied statement; the identically zero polynomial requires separate clarification. For a prime \(p\), write \(\Q_p\) for the field of \(p\)-adic numbers and \(\Z_p\) for its ring of integers. Normalize the absolute value by \(|p|_p=p^{-1}\) and additive Haar measure by \(\operatorname{vol}(\Z_p^n)=1\). Under the standard unweighted interpretation, define
\[
 Z_p(f;s)=\int_{\Z_p^n}|f(x)|_p^s\,\dd x,
 \qquad \operatorname{Re}(s)>0,
\]
and use its meromorphic continuation to define its poles. At zeros of \(f\), the integrand in this initial half-plane is taken to be zero. The source does not specify the integration domain, a test function, or a character; the displayed integral is an explicit interpretive convention.

View \(f\) also as a polynomial map \(\C^n\to\C\). For a nonconstant \(f\) and a point \(a\in f^{-1}(0)\), let \(F_a\) be a local Milnor fibre of the germ of \(f\) at \(a\). Let \(T_{a,j}\) denote the monodromy operator on ordinary cohomology \(H^j(F_a;\C)\), associated with a positively oriented loop around zero. A local monodromy eigenvalue means an eigenvalue of some \(T_{a,j}\), with \(j\geq 0\).

\textbf{Literal question supplied by the record.} For every such \(n\) and \(f\), does there exist a finite set \(S_f\) of primes, which may include the primes dividing denominators of coefficients of \(f\), such that
\[
 \forall p\notin S_f\quad
 \forall s_0\in\C\text{ that is a pole of }Z_p(f;s),\quad
 \exists a\in f^{-1}(0)\ \exists j\geq 0:
 \quad \exp(2\pi i s_0)\in\operatorname{Spec}(T_{a,j})?
\]
The point \(a\) and degree \(j\) may depend on \(p\) and \(s_0\). For a nonzero constant polynomial the displayed integral has no poles, so this convention gives a vacuous condition.

\textbf{Unresolved formulation issue.} The source explicitly uses \(\exp(2\pi i s_0)\) for every complex pole. The usual formulation associated with the title instead uses
\[
 \exp\bigl(2\pi i\operatorname{Re}(s_0)\bigr).
\]
These expressions differ for nonreal poles. The record does not restrict \(s_0\) to real poles or specify that it denotes a real part. Consequently the literal question above and this possible intended formulation must remain distinct pending source review; no correction is silently adopted.
\subsection{Short English statement}
The supplied question asks whether, outside finitely many primes depending on a rational polynomial, each pole of its \(p\)-adic Igusa zeta function determines a local complex monodromy eigenvalue somewhere on its zero set. Its use of the full complex pole instead of its real part leaves the intended formulation unresolved. The supplied partially solved label is dated curation, not an independently verified conclusion.
\subsection{Sources}
\begin{itemize}
\item Local registry: supplied UnsolvedMath record in \texttt{data/unsolvedmath/problems.json}, permanent ID \(30001349\), title \emph{Igusa Monodromy Conjecture}, exact source code \texttt{catalogue entry}, in the Oberwolfach Reports Open Problems collection. No individual proposer is supplied.
\item Supplied report citation: \emph{Singularities (2010)}, Oberwolfach Reports, pp.\ 2405--2470, DOI \texttt{10.4171/owr/2009/43}; \url{https://doi.org/10.4171/owr/2009/43}. This is also the supplied literature link labeled ``original Oberwolfach report.'' The record gives proposed year \(2009\), while the report title as supplied contains \(2010\); these details are retained without reconciliation.
\item Supplied literature link labeled ``introduction to the monodromy conjecture'': \url{https://arxiv.org/abs/2403.03343}.
\item Supplied literature link labeled ``JEMS monodromy results'': \url{https://ems.press/content/serial-article-files/32760}.
\item Literature assessment checked at \texttt{2026-08-21}: the supplied curation labels the problem partially solved and states that the monodromy conjecture is proved in dimension two and many special higher-dimensional cases, but remains open in general. This is supplied dated curation derived from the Oberwolfach Reports material, not independently verified here; its time-sensitive status and cited claims require independent verification, including identification of the precise conjecture variant covered.
\item The registry labels the statement \texttt{corrected\_verified}, explaining this as meaning-preserving editorial and LaTeX cleanup with no source-dependent content added. This supplied editorial label does not establish mathematical correctness or resolve the formulation issues above. The references have not been independently consulted for this draft.
\end{itemize}
\end{document}
